% =====================================================================
% overleaf/02_related_work.tex — LaTeX rendering of en/02_related_work.md v0.17 (10 September 2026)
% Dernière mise à jour : 2026-09-10
% Chapter 2 of the paper, a \section of its own, to be inserted right after overleaf/01_introduction.tex.
% Labels sec:agents and sec:alignment are referenced from the introduction; keep them.
% Ready to paste into the Overleaf project (acmart + natbib: \citep, \citet).
% Derived file: regenerate after every validated pass of en/01_introduction.md.
% Citation keys are those of docs/paper/sources/references.bib (checklist: article/CITATIONS.md).
% Preamble additions:
%   \newcommand{\todo}[1]{\textbf{[#1]}}   % placeholder figures, filled from methode/experience_plan/experiments.yaml
% Vocabulary: Level 0--3 = our ablation levels; exploratory / robust / transferable = SILICA certification steps.
% Comments "% source:" give the repository file a plain-text figure is recomputed from.
% =====================================================================
% =====================================================================

\section{Related work}
\label{sec:related}

\subsection{From random utility to bounded rationality}

Transport engineering has modelled mode choice for fifty years with discrete choice models grounded in random utility maximisation \citep{mcfadden1974conditional,benakiva1985discrete}. Multinomial and mixed logit models express the probability of each mode as a function of travel time, cost and socioeconomic attributes, and are estimated on travel surveys \citep{train2009discrete}. Their strength is statistical: they are calibrated on the population they describe, and their aggregate predictions inherit that calibration. Their weakness is behavioural: taste heterogeneity is represented by static distributions, and the decision rule is a maximisation.

Behavioural research documents departures from this rule. In an online survey of 650 respondents, \citet{adam2025survey} find that habitual car users tend to underestimate the price of the car, and that users accustomed to one mode perceive the others through that habit; they build an agent-based mode-choice model in which such perception filters modulate an otherwise multi-criteria decision. These effects are real but conditional on habit and on the individual, and household travel surveys do not record them. Tabular models therefore cannot represent them, and LLM agents are often proposed as the way to fill this gap.

\subsection{Generative agents in mobility simulation}
\label{sec:agents}

\citet{liu2024toward} propose a conceptual framework in which LLM agents with profiles, memory, perception and decision modules act as proxies for travellers in activity-based demand models. They support the vision with a small proof of concept and identify hybrid modelling, LLM agents combined with established components, as the realistic near-term integration path. CitySim \citep{bougie2025citysim} scales LLM-driven agents to a city and validates them on time use against the 2021 Japanese national time-use survey, reported as visual agreement of daily shares, and on travel volumes and well-being against datasets its authors describe as proprietary, crowd density being compared with estimates derived from smartphone location data; the authors list the non-public data as a limit to reproducibility. Its architecture includes a vehicle-selection module, but no experiment confronts a simulated modal split with an observed one. \citet{chopra2024limits} quantify the cost of giving every agent of a large population its own LLM call and propose archetypes as a remedy, which places inference cost at the centre of any population-scale design.

\citet{lammer2026gta} is the closest published system on the evaluation side. It builds census-aligned personas from the German national travel survey by truncate-replicate-sample, has an LLM write a persona description and a full-day activity plan, then choose one mode per trip among candidate itineraries produced by SUMO and OpenTripPlanner, and closes the loop with a dynamic user equilibrium in SUMO checked against official traffic counts. On a 1\,\% sample of Berlin, 35{,}769 agents, it reports a root-mean-square error of 4.07 on the modal split, 6.07 on trip duration and 6.12 on trip length. This is the first evaluation we know of that confronts an LLM agent population with an observed modal split, and it is the reason our own contribution is a contract rather than a measurement. GTA gives itself no modelled reference: no discrete choice model, no supervised classifier, and no naive floor. Its comparator is geographical, the empirical indicators of the fifteen other German states, and its own Table~2 shows the price of that choice, since predicting Berlin's modal split by simply copying Hamburg's observed one yields an error of 1.99, less than half the 4.07 of the simulation. Its ablations vary the model and the prompt, from 5.42 to 8.23 in aggregate error, but never step outside the LLM. Each agent returns a single mode rather than a distribution, so no distributional metric is available; days are independent, and the authors name multi-day memory and adaptation as future work.

\citet{alves2026evaluating} stand closest on the platform side. They couple GAMA to an external LLM module that decides, on each disruption event, whether an agent should replan its route; they add a persistent memory and compare rule-based and LLM-assisted agents across road-blockage scenarios and population sizes. They report greater adaptability and contextual awareness of LLM agents. Their evaluation, however, has no human ground truth: adaptability is judged against the rule-based baseline, not against observed behaviour, and no distributional metric is computed. Our paper takes the position that neither work occupies. We keep the platform-level design, GAMA plus routing plus LLM plus memory; we adopt from GTA the demand that a simulated modal split be confronted with an observed one; and we add what neither provides: a floor, a calibrated ceiling, informational parity between them, and a distributional reading of the agent's answer.

\subsection{Distributional alignment of LLM populations}
\label{sec:alignment}

A separate literature asks whether language models can stand in for human samples. \citet{argyle2023outofone} showed that a model conditioned on sociodemographic attributes reproduces some response patterns of the corresponding subpopulations, a practice since known as silicon sampling. \citet{meister2024benchmarking} built a benchmark of distributional alignment and found that the way a model expresses a distribution matters as much as the model: asked to verbalise a distribution over options, models align better than when they are sampled, and a model may know a distribution it cannot sample from. Two results then bear directly on our question. \citet{kambhatla2025improving} show that simple supervision improves the alignment of generated distributions across population groups. \citet{huang2025distribution} show that aligning the shift between a source and a target distribution helps models simulate survey response distributions. \citet{nguyen2025prompt} take a population view and select a set of agents, each conditioned on human demonstrations, whose union covers the diversity of a target population.

At the scale of interacting populations, \citet{baronchelli2025emergent} show that decentralised LLM populations form conventions and develop collective biases absent in individual agents, and they state that their work does not treat LLMs as proxies for human participants. \citet{baronchelli2026groupsize} show that group size changes these dynamics in a non-linear, model-dependent way. The SILICA instrument \citep{bintareaf2026silica} turns these concerns into a certification procedure with an explicit perturbation library, option order included; we adopt it as a reading grid, and Section~\ref{sec:gap} states the concern it makes operational.
